\documentclass[fleqn,letterpaper, 12pt]{article}

%%%%%%%%%%%%%%%%%%%%%%%%%%%%%%%%%%%%%%%%%%%%%%%%%%%%
%%%%%%%%%%%%%%%%%   PACKAGES   %%%%%%%%%%%%%%%%%%%%%
%%%%%%%%%%%%%%%%%%%%%%%%%%%%%%%%%%%%%%%%%%%%%%%%%%%%

\usepackage[top = 1.6cm, left = 2cm, right = 2cm ]{geometry}
\usepackage[pdftex]{graphicx}
\usepackage[numbers]{natbib}
\usepackage[french]{babel}
\usepackage{soulutf8}
\usepackage{amsmath}
\usepackage{amssymb}
\usepackage{tikz}
\usepackage[utf8]{inputenc}
\usepackage{longtable}
\usepackage[T1]{fontenc}
\usepackage{epigraph}
\usepackage{fancyhdr}
\usepackage{float}
\usepackage{subfig}
\usepackage{xcolor}
\usepackage{eurosym}
\usepackage{calc}
\usepackage{hyperref}
\usepackage{multirow}
\usepackage{caption}
\usepackage{amsmath}
\usepackage[Algorithme]{algorithm}
\usepackage{algorithmic}
\usepackage{enumerate}
\usepackage{mathtools}
\usepackage{verbatim}
\usepackage{etoolbox}

%%%%%%%%%%%%%%%%%%%%%%%%%%%%%%%%%%%%%%%%%%%%%%%%%%%%
%%%%%%%%%%%%%%%%%   TEMPLATE   %%%%%%%%%%%%%%%%%%%%%
%%%%%%%%%%%%%%%%%%%%%%%%%%%%%%%%%%%%%%%%%%%%%%%%%%%%




\def\changemargin#1#2{\list{}{\rightmargin#2\leftmargin#1}\item[]}
\let\endchangemargin=\endlist 
\newcommand{\newlinealinea}{
	~\\ \hspace*{0.5cm}}
\newcommand{\alinea}{
	\hspace*{0.5cm}}
\newcommand{\alinealong}{
	\hspace*{1.1cm}}
\newcommand{\alignparagraph}{
	\hspace*{0.6cm}}
\newcommand{\red}[1]{
	\textcolor{red}{#1}}
\newcommand{\green}[1]{
	\textcolor{green}{#1}}
\newcommand{\point}{$\bullet\ $}
\makeatletter
\newcommand*{\whiten}[1]{\llap{\textcolor{white}{{\the\SOUL@token}}\hspace{#1pt}}}
\newcommand{\myul}[1]{
	\underline{\smash{#1}}
}
\makeatother
\setlength{\fboxsep}{2pt}
\DeclareMathOperator*{\argmax}{\arg\!\max}
\makeatletter
\@addtoreset{section}{part}
\makeatother  
\renewcommand*\sfdefault{phv}
\renewcommand*\rmdefault{ppl}
\renewcommand\epigraphflush{flushright}
\renewcommand\epigraphsize{\normalsize}
\setlength\epigraphwidth{0.7\textwidth}
\definecolor{titlepagecolor}{cmyk}{0.24,0.92,0.78,0.25}
\definecolor{red}{cmyk}{0, 0.91, 0.91, 0.20}
\DeclareFixedFont{\titlefont}{T1}{phv}{\seriesdefault}{n}{0.375in}
\DeclareFixedFont{\subtitlefont}{T1}{phv}{\seriesdefault}{n}{0.25in}
\cfoot{\thepage}
\pagestyle{fancy}


%%%%%%%%%%%%%%%%%%%%%%%%%%%%%%%%%%%%%%%%%%%%%%%%%%%%
%%%%%%%%%%%%%%%%%  TITLE PAGE %%%%%%%%%%%%%%%%%%%%%%
%%%%%%%%%%%%%%%%%%%%%%%%%%%%%%%%%%%%%%%%%%%%%%%%%%%%


\newcommand\titlepagedecoration{%
	\begin{tikzpicture}[remember picture,overlay,shorten >= -10pt]
	
	\coordinate (aux1) at ([yshift=-70pt]current page.north east);
	\coordinate (aux2) at ([yshift=-460pt]current page.north east);
	\coordinate (aux3) at ([xshift=-6cm]current page.north east);
	\coordinate (aux4) at ([yshift=-150pt]current page.north east);
	
	\begin{scope}[titlepagecolor!40,line width=12pt,rounded corners=12pt]
	\draw
	(aux1) -- coordinate (a)
	++(225:5) --
	++(-45:5.1) coordinate (b);
	\draw[shorten <= -10pt]
	(aux3) --
	(a) --
	(aux1);
	\draw[opacity=0.6,titlepagecolor,shorten <= -10pt]
	(b) --
	++(225:2.2) --
	++(-45:2.2);
	\end{scope}
	\draw[titlepagecolor,line width=8pt,rounded corners=8pt,shorten <= -10pt]
	(aux4) --
	++(225:0.8) --
	++(-45:0.8);
	\begin{scope}[titlepagecolor!70,line width=6pt,rounded corners=8pt]
	\draw[shorten <= -10pt]
	(aux2) --
	++(225:3) coordinate[pos=0.45] (c) --
	++(-45:3.1);
	\draw
	(aux2) --
	(c) --
	++(135:2.5) --
	++(45:2.5) --
	++(-45:2.5) coordinate[pos=0.3] (d);   
	\draw 
	(d) -- +(45:1);
	\end{scope}
	\end{tikzpicture}%
}

%%%%%%%%%%%%%%%%%%%%%%%%%%%%%%%%%%%%%%%%%%%%%%%%%%%%
%%%%%%%%%%%%%%%%%%  ALGORITHM  %%%%%%%%%%%%%%%%%%%%%
%%%%%%%%%%%%%%%%%%%%%%%%%%%%%%%%%%%%%%%%%%%%%%%%%%%%

\renewcommand{\algorithmicrequire}{\textbf{Requis:}}
\renewcommand{\algorithmicensure}{\textbf{Promet:}}
\renewcommand{\algorithmicend}{\textbf{fin}}
\renewcommand{\algorithmicif}{\textbf{si}}
\renewcommand{\algorithmicthen}{\textbf{alors}}
\renewcommand{\algorithmicelse}{\textbf{sinon}}
\renewcommand{\algorithmicelsif}{\algorithmicelse\ \algorithmicif}
\renewcommand{\algorithmicendif}{\algorithmicend\ \algorithmicif}
\renewcommand{\algorithmicfor}{\textbf{pour}}
\renewcommand{\algorithmicforall}{\textbf{pour chaque}}
\renewcommand{\algorithmicdo}{\textbf{faire}}
\renewcommand{\algorithmicendfor}{\algorithmicend\ \algorithmicfor}
\renewcommand{\algorithmicwhile}{\textbf{tant que}}
\renewcommand{\algorithmicendwhile}{\algorithmicend\ \algorithmicwhile}
\renewcommand{\algorithmicloop}{\textbf{boucle}}
\renewcommand{\algorithmicendloop}{\algorithmicend\ \algorithmicloop}
\renewcommand{\algorithmicrepeat}{\textbf{répéter}}
\renewcommand{\algorithmicuntil}{\textbf{jusqu'à}}
\renewcommand{\algorithmicprint}{\textbf{afficher}}
\renewcommand{\algorithmicreturn}{\textbf{retourner}}
\renewcommand{\algorithmictrue}{\textbf{vrai}}
\renewcommand{\algorithmicfalse}{\textbf{faux}}

\renewcommand{\algorithmicand}{\textbf{et}}
\renewcommand{\algorithmicor}{\textbf{ou}}

%%%%%%%%%%%%%%%%%%%%%%%%%%%%%%%%%%%%%%%%%%%%%%%%%%%%
%%%%%%%%%%%%%%%%%   CUSTOM TIKZ    %%%%%%%%%%%%%%%%%
%%%%%%%%%%%%%%%%%%%%%%%%%%%%%%%%%%%%%%%%%%%%%%%%%%%%

\newcommand{\tikzmultiline}[6]{
	% {from}{to}{arrows}{offset from}{offset to}{arrow angle}
	\pgfmathsetmacro{\pgfminus}{#3-1}%
	\foreach \i in {0,...,\pgfminus}{
		 \pgfmathsetmacro{\pgfresult}{\i*100/\pgfminus}%
		\draw[-] (#1.{(\i)*#6+#4 - \pgfminus*#6/2}) -- (#2.{-(\i)*#6+#5  + \pgfminus*#6/2});}
}



%%%%%%%%%%%%%%%%%%%%%%%%%%%%%%%%%%%%%%%%%%%%%%%%%%%%
%%%%%%%%%%%%%%%%   NOTATIONS    %%%%%%%%%%%%%%%%%%%%
%%%%%%%%%%%%%%%%%%%%%%%%%%%%%%%%%%%%%%%%%%%%%%%%%%%%







%%%%%%%%%%%%%%%%%%%%%%%%%%%%%%%%%%%%%%%%%%%%%%%%%%%%
%%%%%%%%%%%%%%%   INFORMATIONS  %%%%%%%%%%%%%%%%%%%%
%%%%%%%%%%%%%%%%%%%%%%%%%%%%%%%%%%%%%%%%%%%%%%%%%%%%

\renewcommand{\title}{\titlefont Philosophie des sciences \\ \subtitlefont Synthèse \par}
\newcommand{\subject}{Philosophie des sciences}
\newcommand{\student}{Benjamin André}
\renewcommand{\date}{\today}


%%%%%%%%%%%%%%%%%%%%%%%%%%%%%%%%%%%%%%%%%%%%%%%%%%%%

\begin{document}
	\lhead{\student}
	\rhead{\subject}
	\pagenumbering{gobble}
	\begin{titlepage}
		%
		\noindent
		%
		\newgeometry{bottom = 2cm, top = 2.5cm}
		\begin{center}
			\includegraphics[scale=1.2]{Images/UMONS}\\
			\vspace*{0.3cm}
			\includegraphics[scale=0.23]{Images/FS_Logo}\\
			\vspace*{2.5cm}
			%
			\title
			\vspace*{2.5cm}
			\date
			%
		\end{center}
		\vspace*{3cm}
		\hfill
		%
		\begin{minipage}{0.18\linewidth}
			\begin{flushright}
				\rule{0.5pt}{50pt}
			\end{flushright}
		\end{minipage}
		%
		\begin{minipage}{0.8\linewidth}
			\begin{flushleft}
				\textsf{\textbf{Sujet :}} \subject\\
				\textsf{\textbf{Étudiant:}} \student
			\end{flushleft}
		\end{minipage}
		%
		\vspace*{\fill}                                                             
		%
		\begin{center}
			Faculté des Sciences $\bullet$ Université de Mons $\bullet$ 
			Place du Parc 20 $\bullet$ B-7000 Mons
		\end{center}
		\titlepagedecoration
	\end{titlepage}
	\newgeometry{top = 3cm, left = 2cm, right = 2cm, bottom=2.5cm}
	\setcounter{tocdepth}{1}\tableofcontents\newpage
	\pagenumbering{arabic}\setcounter{page}{1}
	
	
	
	
	
	
	
	\section{Introduction}
	
	Buts : 
	\begin{enumerate}
		\item Connaître les principales conceptions épistémologiques
		\item Savoir argumenter
		\item Réfléchir $\leftrightarrow$  donner son opinion
	\end{enumerate}
	
	Oral :
	\begin{itemize}
		\item Une question de compréhension / réflexion
		\item Une question de connaissance générale
	\end{itemize}

	\section{Qu'est-ce que la philosophie ?}
	
	Pratique réflexive (et souvent auto-réflexive) et argumentative, qui aborde de manière critique les divers domaines des connaissances et des croyances en privilégiant le doute et le questionnement sur les réponses.
	
	\section{L'inductivisme}
		
	\begin{itemize}
		\item Francis Bacon (1561-1626)
		\item Observation $\rightarrow$ Expérimentation $\rightarrow$ Théorie scientifique vraie
		\item Le nombre d'énoncés d'observation doit être \textbf{élevé}
		\item Les observations doivent être répétées dans une grande \textbf{variété de conditions}
		\item \textbf{Aucun} énoncé d'observation ne peut entrer en contradiction avec la loi.
	\end{itemize}
	
	
	\paragraph{Critiques}
	\begin{itemize}
		\item Ce n'est logiquement pas valide (de dire que si c'est vrai pour un grand nombre d'observations, c'est vrai toujours)
		\item Quels sont les critères de variations ?
		\item Combien d'observations sont nécessaires ?
		\item L'observation est une action plus qu'une pure perception
		\item Les énoncés d'observation supposent une théorie
		\item Les énoncés d'observation sont faillibles
	\end{itemize}
	
	
	\section{Descartes}
	
	La philosophie c'est l'étude de la sagesse et de ses principes. Il n'y a que Dieu qui soit parfaitement sage, qui connaît la vérité de toutes choses. Il n'y a personne qui ne souhaite quelque chose de plus important que les biens matériels : la vérité.
	Différents degrés de sagesse
	\begin{enumerate}
		\item Les notions si claires qu'on peut les acquérir sans méditation
		\item Tout ce que l'expérience des sens fait connaître
		\item Ce que la conversation des autres hommes nous enseigne
		\item Lecture de livre écrit par des hommes sages
	\end{enumerate}
	En plus de cela, il y a la révélation divine, ne fonctionnant pas par degré car amenant à une créance infaillible.
	
	Il y aurait un 5e degré, qui consisterait à chercher les causes et principes de tout ce qu'on peut savoir. Les gens qui les cherchent sont des philosophes.
	
	Il faut commencer par se former une morale, et conduire sa raison pour découvrir les vérités qu'on ignore. Ensuite, on peut se mettre à la vraie philosophie, en commençant par la métaphysique, qui contient l'explication des principaux attributs de Dieu, de l'immatérialité de nos âmes et de toutes les notions claires et simples qui sont en nous. Ensuite, la physique. Comment l'univers est composé, puis la nature et les corps qui s'y trouvent (~éléments). En suite il faut étudier en particulier la nature des plantes, animaux et de l'homme, pour trouver les autres sciences qui seraient utiles.
	
	La philosophie est comme un arbre, dont les racines sont la métaphysique et le tronc la physique. Cependant c'est au bout des branches (médecine, mécanique et morale) qu'on cueille les fruits de la philosophie.
	
	
	\section{Le falsificationisme}
	
	\begin{itemize}
		\item Karl Popper (1902-1994)
		\item "La science vient de théories qu'il faut rigoureusement soumettre à l'observation et à l'expérimentation".
		\item Si on ne peut pas imaginer le contraire d'un énoncé, il n'est pas falsifiable et donc pas scientifique.
	\end{itemize}
	
	\paragraph{Le falsificationisme sophistiqué}
	\begin{itemize}
		\item Degré de falsifiabilité (quelque chose qui serait plus facile a falsifié mais qui ne l'a pas été). Parmi plusieurs théories concurrentes, préférer celle qui prédit un type de phénomènes supplémentaires.
		\item Théories ad hoc (la théorie modifiée ne peut être testé sans que le test ne touche également la théorie de départ et la théorie modifiée est moins falsifiable que celle qui l'a précédée) et non ad hoc 
		\item La confirmation (si une théorie échoue a un test, elle est falsifiée. Sinon, elle est confirmée). Il y a plus de progrès quand une théorie prudente est falsifiée et une théorie audacieuse confirmée.
	\end{itemize}
	
	\paragraph{Critiques}
	\begin{itemize}
		\item Une théorie scientifique ne peut jamais être vraie. (Tous les corbeaux sont noirs)
		\item Dépendance de l'observation par rapport à la théorie
		\item Faillibilité des falsifications
		\item Complexité des situations de tests réalistes
		\item Raisons historiques
	\end{itemize}
	
	\section{Lakatos}
	
	\begin{itemize}
		\item Imre Lakatos (1922-1974)
		\item Noyau dur (hypothèses très générales infalsifiables. Toute inadéquation doit se travailler pour tomber sur ce noyau dur. Exemple : les épicycles de Copernic)
		\item Heuristique négative : le noyau dur d'une théorie ne doit être modifié au rejeté
		\item Heuristique positive : lignes directrices d'un programme de recherche. Hypothèses venant englober le noyau dur. Modification non ad hoc autorisée.
	\end{itemize}

	On protège donc plus les énoncés universels qu'avec Popper. De plus, les confirmations sont plus importantes. La science progresse grâce aux divers programmes de recherche.

	Une théorie est scientifique $\Leftrightarrow$ :
	\begin{itemize}
		\item Elle est falsifiable
		\item Elle a une cohérence permettant d'inclure la recherche future
		\item Elle mène à la découverte de nouveaux phénomènes
	\end{itemize}
	
	\paragraph{Critiques}
	\begin{itemize}
		\item Une théorie n'est jamais fausse
		\item Difficulté de comparaison des programmes de recherche (plus facile pour un historien qu'un décideur)
		\item Dépendance des programmes de recherche
		\item Vision économiste
		\item Importance de la physique
		\item Moins grande falsifiabilité
		\item Limites floues au noyau dur
	\end{itemize}

	\section{Kuhn}
	
	\begin{itemize}
		\item Thomas Samuel Kuhn (1922-1996)
		\item Pré-science (désaccord sur les fondements de la science) $\rightarrow$ Science normale (paradigme)$\rightarrow$ Crise et révolution (on découvre quelque chose de complètement nouveau / on a plus confiance en notre capacité de résoudre les énigmes de la nature) $\rightarrow$ Nouvelle science normale (nouveau paradigme)
		\item Prend en compte le caractère sociologique des communautés scientifiques, ce qui le différencie de Popper et Lakatos.
	\end{itemize}

	\paragraph{Critiques}
	\begin{itemize}
		\item Relativisme
		\item Manque de critère rationnel pour justifier le changement de paradigme
		\item La science ne progresse pas vers la vérité
		\item Présupposé : la science est supérieure à d'autres formes de savoir
	\end{itemize}
	
	
	\section{Feyerabend}
	
	\begin{itemize}
		\item Paul Feyerabend (1924-1994)
		\item "Tout est bon" (on a fait des avancées scientifique sous de nombreuses méthodologies différentes)
		\item L'incommensurabilité (certaines théories sont incomparables, comme la relativité générale et la mécanique quantique. Après, il ne reste que notre jugement subjectif)
		\item La science n'est pas nécessairement supérieure aux autres disciplines (la science étant elle même incommensurable, on ne peut pas rejeter le marxisme car infalsifiable par exemple)
		\item La liberté de l'individu
		\item Une organisation fixée et universelle serait préjudiciable à la science car elle néglige les conditions physiques et historiques. Elle rendrait la science moins facilement adaptable et plus dogmatique (nldr : risque de tout perdre lors d'un changement historique ?)
		\item Lorsqu'on compare deux théories, après comparaison des conséquences et prédictions, ce qu'il reste, ce sont nos désirs subjectifs.
		\item Feyerabend a écrit "Contre la méthode : Esquisse d'une théorie anarchique de la science"
	\end{itemize}
	
	
	
	\section{Deleuze}
	
	Ne sert ni l'\'Etat ni l'\'Eglise, mais à attrister. Faire de la bêtise quelque chose de honteux. Dénoncer toutes les fictions pour que les hommes soient libres. Entreprise de démystification. Les philosophes aujourd'hui sont devenus des comètes (Nietzsche), qu'il nous faut regarder. Opposer l'homme libre aux valeurs et puissances établies.
	
	Le philosophe est un ami de la sagesse, de la vérité face à l'épreuve du sens et de la valeur. Les philosophes de religion, d'état sont des traitres.
	
	Le philosophe est toujours contre son temps, critique du monde actuel, forme des modèles qui ne sont ni éternels ni historiques, en opposition à la bêtise qui est éternelle. En proposant un modèle intempestif, il y a des vérités plus durables que les vérités historiques et éternelles réunies : les vérités du temps à venir.
	
	La chaîne des philosophes est une chaîne brisée. Il n'y a pas de philosophie éternelle ni historique. La philosophie toujours intempestive, intempestive à chaque époque.
	
	
	
	


	
	
	
	
\end{document}